\section*{Capítulo \ref{ch:g5:weather-and-climate} --- Tempo e clima}

\begin{solution}{exo:g5:weather-and-climate:1}
(a) \omterm{def:g1:watching-the-weather:weather}{Tempo} (só hoje de manhã). (b) \omterm{def:g5:weather-and-climate:climate}{Clima} (todo verão). (c) \omterm{def:g1:watching-the-weather:weather}{Tempo} (ontem à
tarde). (d) \omterm{def:g5:weather-and-climate:climate}{Clima} (todos os meses, ano após ano).
\end{solution}

\begin{solution}{exo:g5:weather-and-climate:2}
O \omterm{def:g5:weather-and-climate:climate}{clima} de um lugar é o \omterm{def:g1:watching-the-weather:weather}{tempo} que ele costuma ter, mês a mês, ao longo de
muitos anos. Os meteorologistas fazem a média das medições de 30 anos.
\end{solution}

\begin{solution}{exo:g5:weather-and-climate:3}
O mês mais frio é dezembro (cerca de 6~$^\circ$C; janeiro é quase tão
frio). O mês mais chuvoso é novembro (cerca de 160~mm). Não: a curva nunca
fica abaixo de 0~$^\circ$C.
\end{solution}

\begin{solution}{exo:g5:weather-and-climate:4}
Julho é o mês mais quente, com cerca de 5~$^\circ$C. Nenhum mês passa de
10~$^\circ$C.
\end{solution}

\begin{solution}{exo:g5:weather-and-climate:5}
O \omterm{def:g3:continents-and-oceans:equator}{equador}: a zona quente. O \omterm{def:g3:continents-and-oceans:poles}{Polo Norte}: a zona fria (a do norte). A meio
caminho entre o \omterm{def:g3:continents-and-oceans:equator}{equador} e o \omterm{def:g3:continents-and-oceans:poles}{Polo Sul}: a zona temperada do sul (ela fica
entre o \omterm{def:g5:weather-and-climate:tropics}{trópico} de Capricórnio e o Círculo Polar Antártico).
\end{solution}

\begin{solution}{exo:g5:weather-and-climate:6}
Hilo: a zona quente. Phoenix e Seattle: a zona temperada do norte.
Utqiaġvik: a zona fria do norte, além do Círculo Polar Ártico.
\end{solution}

\begin{solution}{exo:g5:weather-and-climate:7}
Ano inteiro: $147 + 96 + 106 + 81 + 48 + 37 + 15 + 25 + 41 + 99 + 160 + 145
= 1\,000$ milímetros. Dezembro, janeiro e fevereiro:
$145 + 147 + 96 = 388$ milímetros, mais de um terço da \omterm{def:g4:clouds-rain-snow:precipitation}{chuva} do ano em
três meses.
\end{solution}

\begin{solution}{exo:g5:weather-and-climate:8}
Phoenix: $35 - 14 = 21$ graus. Hilo: $25 - 22 = 3$ graus. As estações
mudam muito mais a temperatura em Phoenix; em Hilo, ela é quase a mesma o
ano todo.
\end{solution}

\begin{solution}{exo:g5:weather-and-climate:9}
Um \omterm{def:g5:weather-and-climate:climate}{clima} de \omterm{def:g5:weather-and-climate:rainforest}{floresta tropical úmida}: faz calor o ano todo (27 ou
28~$^\circ$C, sem inverno), e cai muita \omterm{def:g4:clouds-rain-snow:precipitation}{chuva} em todos os meses (pelo menos
160~mm, e mais de 2\,000~mm no ano). A cidade fica, mais provavelmente, na
zona quente, perto do \omterm{def:g3:continents-and-oceans:equator}{equador}.
\end{solution}

\begin{solution}{exo:g5:weather-and-climate:10}
Primeira explicação: a praia fica na zona quente, onde faz calor o ano
todo, e o jardim com \omterm{def:g4:clouds-rain-snow:precipitation}{neve} numa zona temperada ou fria, onde janeiro cai no
inverno. Segunda explicação: a praia fica na metade sul da Terra, onde
janeiro cai no verão, e o jardim com \omterm{def:g4:clouds-rain-snow:precipitation}{neve} na metade norte, onde janeiro
cai no inverno: as estações são invertidas entre as duas metades do \omterm{def:g3:continents-and-oceans:globe}{globo}.
\end{solution}

\begin{solution}{exo:g5:weather-and-climate:11}
Temperaturas: Phoenix é muito mais quente, cerca de 35~$^\circ$C em julho
contra cerca de 20~$^\circ$C em agosto em Seattle, e o seu janeiro (cerca
de 14~$^\circ$C) já é tão quente quanto o maio de Seattle. \omterm{def:g4:clouds-rain-snow:precipitation}{Chuva}: menos de
200~mm por ano em Phoenix, um \omterm{def:g5:weather-and-climate:desert}{deserto}, contra cerca de 1\,000~mm em
Seattle. Então elas não têm o mesmo \omterm{def:g5:weather-and-climate:climate}{clima}. As zonas nos dizem, de modo
aproximado, quanto o Sol esquenta um lugar, do \omterm{def:g3:continents-and-oceans:equator}{equador} até os \omterm{def:g3:continents-and-oceans:poles}{polos}; elas
não nos dizem quanta \omterm{def:g4:clouds-rain-snow:precipitation}{chuva} cai, nem o efeito do \omterm{def:g2:water-on-earth:ocean}{mar}, das \omterm{def:g1:landscapes-around-us:landform}{montanhas} ou dos
ventos.
\end{solution}
